% withoutpreface: 跳过承诺书页 + 编号专用页(电子版论文不需此两页)
% 用户纸质提交时手动添加,详见 cls 中 \@mcm@preface 开关
\documentclass[withoutpreface]{cumcmthesis}

% ======================== 包引入 ========================
\usepackage{graphicx}
\usepackage{booktabs}
\usepackage{amsmath,amssymb}
\usepackage{geometry}
\usepackage{enumitem}
\usepackage{multirow}
\usepackage{tabularx}
\usepackage{float}
\usepackage{subcaption}
\usepackage[titles]{tocloft}
\usepackage{hyperref}
\usepackage{xcolor}
\usepackage{listings}
\usepackage{appendix}
\usepackage{longtable}

% ======================== 目录紧凑设置 ========================
\setcounter{tocdepth}{2}  % 显示到 subsection
\setlength{\cftbeforesecskip}{0pt plus 0pt}
\setlength{\cftbeforesubsecskip}{0pt plus 0pt}
\renewcommand{\cftdotsep}{2}

% ======================== 页面设置 ========================
\geometry{a4paper, left=2.5cm, right=2.5cm, top=2.5cm, bottom=2.5cm}

% 摘要左右边距与正文统一：去掉默认 \quotation 造成的左右内缩（约 2.5em），
% 使摘要正文占满全文宽度，与正文段落左右边距一致。
\renewenvironment{abstract}{%
  \small
  \begin{center}%
    {\bfseries\abstractname\vspace{-.5em}}%
  \end{center}%
  \par\noindent\ignorespaces
}{%
  \par
}
\setlength{\parskip}{2pt plus 0.5pt minus 0.5pt}
\linespread{1.0}
\setlength{\textfloatsep}{8pt plus 2pt minus 2pt}
\setlength{\intextsep}{8pt plus 2pt minus 2pt}

% ======================== 图表搜索路径 ========================
\graphicspath{{../figures/}{figures/}}

% ======================== 标题信息 ========================
\title{全国大学生数学建模竞赛论文}
\author{作者1 \and 作者2 \and 作者3}
\date{\today}

% ======================== 正文 ========================
\begin{document}

\maketitle

\begin{abstract}
在此填写摘要内容。摘要应包含问题概述、模型方法、主要结果和结论。
\par\textbf{关键词：} 关键词1；关键词2；关键词3
\end{abstract}

\tableofcontents
\newpage

% ======================== 问题重述 ========================
\section{问题重述}
\subsection{问题背景}
在此简述问题背景。

\subsection{问题提出}
\begin{enumerate}[label=(\arabic*)]
    \item 问题一的描述。
    \item 问题二的描述。
    \item 问题三的描述。
\end{enumerate}

% ======================== 模型假设 ========================
\section{模型假设}
\begin{enumerate}[label=\textbf{假设\arabic*：}]
    \item 假设内容一。
    \item 假设内容二。
    \item 假设内容三。
\end{enumerate}

% ======================== 符号说明 ========================
\section{符号说明}
\begin{table}[H]
\centering
\begin{tabular}{lll}
\toprule
\textbf{符号} & \textbf{含义} & \textbf{单位} \\
\midrule
$x$ & 示例符号 & 示例单位 \\
\bottomrule
\end{tabular}
\end{table}

% ======================== 问题一 ========================
\section{问题一的分析与求解}
\subsection{问题分析}
在此分析问题一。

\subsection{模型建立}
在此建立模型。

\subsection{模型求解}
在此求解模型。

\subsection{结果}
在此展示结果。

% ======================== 问题二 ========================
\section{问题二的分析与求解}
\subsection{问题分析}
在此分析问题二。

\subsection{模型建立}
在此建立模型。

\subsection{模型求解}
在此求解模型。

\subsection{结果}
在此展示结果。

% ======================== 问题三 ========================
\section{问题三的分析与求解}
\subsection{问题分析}
在此分析问题三。

\subsection{模型建立}
在此建立模型。

\subsection{模型求解}
在此求解模型。

\subsection{结果}
在此展示结果。

% ======================== 灵敏度分析 ========================
\section{灵敏度分析}
对模型参数进行灵敏度分析，验证模型的稳定性。

% ======================== 模型评价与推广 ========================
\section{模型的评价与推广}
\subsection{模型优点}
\begin{enumerate}
    \item 优点一。
    \item 优点二。
\end{enumerate}

\subsection{模型缺点}
\begin{enumerate}
    \item 缺点一。
    \item 缺点二。
\end{enumerate}

\subsection{模型推广}
在此讨论模型的推广方向。

% ======================== 参考文献 ========================
\begin{thebibliography}{99}
\bibitem{ref1} 参考文献1
\bibitem{ref2} 参考文献2
\end{thebibliography}

% ======================== 附录 ========================
\begin{appendices}
\section{代码}
\begin{lstlisting}[language=Python]
# Python代码示例
print("Hello, World!")
\end{lstlisting}
\end{appendices}

\end{document}
